\section{Results and Discussion}
\subsection{Verified Speedup}
Table~\ref{tab:verified} reports the verified results. All methods remained below \texttt{-O3} on average: the upper limit of every 95\% confidence interval is below 1. The proposed method (0.942) and TPE (0.952) are statistically indistinguishable: no paired Wilcoxon test of the proposed method against any other method is significant (Table~\ref{tab:wilcoxon}, all $p > 0.3$). Removing the warm start changes neither the speedup (0.933, $p = 0.75$) nor the conclusion, so the warm start does not explain the results. The median verified speedup is between 0.999 and 1.001 for every method, so the selected sequence of a typical run is approximately as fast as \texttt{-O3}; the verified beat rates (44--64\%) count any speedup above 1 without a noise margin and should be read accordingly.

\subsection{Search-Time versus Verified Results}
Table~\ref{tab:search} lists the search-time measures. Selecting the best of 20 candidates, each timed as the minimum of three runs, favors lucky timings: for example, the search-time beat rate of the proposed method is 64\% against a verified beat rate of 44\%, and its best-of-20 geometric-mean speedup (0.956) falls to 0.942 after re-timing. Search-time numbers are therefore reported only as measures of search behavior, not of final performance.

\subsection{Per-Benchmark Results}
Table~\ref{tab:perbench} shows that no method is best on every benchmark. On 2MM and covariance all methods are close to \texttt{-O3}. The proposed method is best on GEMM (0.983) but has a large variance on Floyd--Warshall (standard deviation 0.205), where TPE (0.980) and Random Search (0.960) are better. With five seeds per cell, the standard deviations are indicative only.

\subsection{Failed Builds and Tuning Cost}
Table~\ref{tab:cost} gives the failure rate and the decomposition of tuning time. Executing candidate programs accounts for 84--87\% of the tuning time of every method, while the remainder (model fitting, failed builds, overhead) is only 3--4\%. The lower tuning time of the proposed method (51.6\,s against 64.7--68.1\,s) therefore comes from the execution component: the candidates it evaluated ran faster on average than those of the baselines, not from a cheaper algorithm. Model-fitting time cannot be separated from time lost in failed builds, because failed builds are not timed individually. The proposed method has the lowest failure rate (2.0\%); the variant without warm start is also low (2.8\%).

\subsection{Statistical Comparison}
Table~\ref{tab:wilcoxon} compares the proposed method with each other method on the 25 matched runs. Table~\ref{tab:explore} reports exploratory comparisons of search cost, chosen after inspecting the means. With 18 comparisons the Bonferroni threshold is $0.05/18 \approx 0.0028$. After this correction, two results remain significant: the lower tuning time of both proposed variants against all three baselines, and the lower number of failed builds of the proposed method compared with TPE. The reduction in evaluations needed to beat \texttt{-O3} against TPE ($p = 0.003$) is borderline, and all other differences are not significant. The seeds on one benchmark are not independent, so the $p$-values are somewhat optimistic.

\subsection{Overall Interpretation}
With five seeds and five benchmarks, no method achieves a verified speedup significantly different from the others, and all remain below \texttt{-O3} on average. This is expected for the setting studied: each sequence replaces the \texttt{-O3} pipeline, only 25 passes in 15 positions are available, and the budget is 20 evaluations. The warm start does not change the verified speedup. The measurable differences concern search behavior: the proposed method wasted fewer evaluations on failed builds than TPE and used less time per search than all baselines because its candidates executed faster. These results do not show that the proposed method finds better final sequences than the baselines, and they should be read as an initial comparison under a small budget, not as evidence of general superiority.

\begin{table*}[t]
\caption{Verified results (25 runs per method; correctness checked, minimum of ten timings). Speedups are relative to \texttt{-O3} (higher is better). Robustness is the standard deviation across seeds averaged over benchmarks; the compile-time ratio is relative to \texttt{-O3}.}
\label{tab:verified}
\centering\small
\begin{tabular}{@{}lccccccc@{}}
\toprule
Method & Geometric-mean speedup [95\% CI] & Median & Beat \texttt{-O3} (\%) & Robustness & Compile-time ratio & Valid (\%)\\
\midrule
Random Search & 0.947 [0.906, 0.982] & 0.999 & 44 & 0.065 & 0.931 & 100\\
Genetic Algorithm & 0.917 [0.860, 0.966] & 1.000 & 48 & 0.095 & 0.901 & 100\\
TPE & 0.952 [0.923, 0.977] & 1.000 & 44 & 0.059 & 0.844 & 100\\
Proposed & 0.942 [0.885, 0.981] & 1.000 & 44 & 0.072 & 0.948 & 100\\
Proposed, no warm start & 0.933 [0.876, 0.981] & 1.001 & 64 & 0.097 & 1.107 & 100\\
\bottomrule
\end{tabular}
\end{table*}

\begin{table}[t]
\caption{Search-time results (minimum of three timings, best of 20 evaluations; biased towards favorable timings)}
\label{tab:search}
\centering\footnotesize
\begin{tabular}{@{}lccc@{}}
\toprule
Method & Best-of-20 speedup & Beat \texttt{-O3} (\%) & Evals.\ to beat \texttt{-O3}\\
\midrule
Random Search & 0.954 & 48 & 12.3\\
Genetic Algorithm & 0.928 & 52 & 12.7\\
TPE & 0.945 & 48 & 14.8\\
Proposed & 0.956 & 64 & 9.9\\
Proposed, no warm start & 0.940 & 68 & 10.0\\
\bottomrule
\end{tabular}
\end{table}

\begin{table*}[t]
\caption{Verified speedup per benchmark (mean $\pm$ standard deviation over five seeds)}
\label{tab:perbench}
\centering\small
\begin{tabular}{@{}lccccc@{}}
\toprule
Benchmark & Random & GA & TPE & Proposed & Proposed, no warm start\\
\midrule
2MM & 0.979 $\pm$ 0.049 & 0.979 $\pm$ 0.052 & 0.980 $\pm$ 0.048 & 0.979 $\pm$ 0.048 & 0.979 $\pm$ 0.049\\
Covariance & 0.998 $\pm$ 0.007 & 1.000 $\pm$ 0.003 & 0.983 $\pm$ 0.039 & 1.000 $\pm$ 0.002 & 1.001 $\pm$ 0.001\\
Floyd--Warshall & 0.960 $\pm$ 0.053 & 0.816 $\pm$ 0.213 & 0.980 $\pm$ 0.047 & 0.863 $\pm$ 0.205 & 0.923 $\pm$ 0.176\\
GEMM & 0.963 $\pm$ 0.052 & 0.919 $\pm$ 0.094 & 0.901 $\pm$ 0.083 & 0.983 $\pm$ 0.039 & 0.878 $\pm$ 0.182\\
Jacobi-2D & 0.860 $\pm$ 0.163 & 0.920 $\pm$ 0.111 & 0.925 $\pm$ 0.078 & 0.926 $\pm$ 0.063 & 0.927 $\pm$ 0.076\\
\bottomrule
\end{tabular}
\end{table*}

\begin{table}[t]
\caption{Failed builds and tuning time per search (seconds)}
\label{tab:cost}
\centering\footnotesize
\setlength{\tabcolsep}{3pt}
\begin{tabular}{@{}lcccccc@{}}
\toprule
Method & Failed/search & Failed (\%) & Total & Compile & Execute & Other\\
\midrule
Random Search & 1.12 & 5.6 & 65.1 & 6.3 & 56.8 & 2.0\\
Genetic Algorithm & 0.64 & 3.2 & 64.7 & 6.7 & 56.0 & 2.0\\
TPE & 1.32 & 6.6 & 68.1 & 6.3 & 59.5 & 2.3\\
Proposed & 0.40 & 2.0 & 51.6 & 6.6 & 43.2 & 1.8\\
Proposed, no warm start & 0.56 & 2.8 & 54.4 & 7.2 & 45.5 & 1.7\\
\bottomrule
\end{tabular}
\end{table}

\begin{table}[t]
\caption{Paired Wilcoxon signed-rank test on verified speedup (25 pairs). Difference is proposed minus other, median.}
\label{tab:wilcoxon}
\centering\small
\begin{tabular}{@{}lcc@{}}
\toprule
Comparison & Median difference & $p$\\
\midrule
Proposed vs Random Search & $+0.0001$ & 0.542\\
Proposed vs Genetic Algorithm & $+0.0004$ & 0.560\\
Proposed vs TPE & $+0.0008$ & 0.325\\
Proposed vs no warm start & $-0.0001$ & 0.751\\
\bottomrule
\end{tabular}
\end{table}

\begin{table*}[t]
\caption{Exploratory paired comparisons of search cost (median difference, first method minus second; $p$ from the Wilcoxon signed-rank test; negative differences favor the first method; Bonferroni threshold 0.0028)}
\label{tab:explore}
\centering\small
\begin{tabular}{@{}lcccccc@{}}
\toprule
 & \multicolumn{2}{c}{Evaluations to beat \texttt{-O3}} & \multicolumn{2}{c}{Failed builds} & \multicolumn{2}{c}{Tuning time (s)}\\
\cmidrule(lr){2-3}\cmidrule(lr){4-5}\cmidrule(lr){6-7}
Comparison & Median diff. & $p$ & Median diff. & $p$ & Median diff. & $p$\\
\midrule
Proposed vs Random Search & 0.00 & 0.135 & 0.00 & 0.006 & $-8.63$ & $<0.001$\\
Proposed vs Genetic Algorithm & 0.00 & 0.162 & 0.00 & 0.516 & $-10.67$ & $<0.001$\\
Proposed vs TPE & $-1.00$ & 0.003 & $-1.00$ & $<0.001$ & $-13.93$ & $<0.001$\\
No warm start vs Random Search & 0.00 & 0.262 & 0.00 & 0.035 & $-12.25$ & $<0.001$\\
No warm start vs Genetic Algorithm & 0.00 & 0.179 & 0.00 & 0.848 & $-10.42$ & $<0.001$\\
No warm start vs TPE & $-1.00$ & 0.006 & 0.00 & 0.018 & $-13.20$ & $<0.001$\\
\bottomrule
\end{tabular}
\end{table*}


\section{Limitations and Future Work}
The study has several limitations. First, only five PolyBench/C benchmarks were evaluated on one machine under WSL2, without CPU pinning or frequency control. The CPU model and frequency-governor state were not recorded in the retained experiment metadata, limiting exact hardware-level reproducibility. Second, the search budget was limited to 20 evaluations per run. Although five random seeds were used, this budget may be insufficient to identify high-performing pass sequences, and larger budgets could change the relative performance of the methods. Third, every candidate sequence replaces the LLVM \texttt{-O3} optimization pipeline, and the search space is limited to 25 optimization passes arranged in 15 positions. Applying sequences after \texttt{-O3}, or exploring a larger pass set, was not evaluated. Fourth, only the proposed method used warm-start sequences, selected from other programs within the same evaluated benchmark set rather than from a disjoint external training set. The no-warm-start ablation did not show a meaningful improvement in verified speedup, so these experiments do not establish that warm-start transfer generalizes to unseen programs; this requires a separate benchmark-disjoint evaluation. Fifth, search-time measurements can favor candidates with unusually low observed execution times, and the search-cost comparisons are exploratory. Seeds on the same benchmark may also be correlated. Sixth, the verified beat rate counts any measured speedup above 1 without accounting for a minimum practically meaningful improvement. Seventh, tuning time could not be fully separated into model-fitting time and time spent on failed builds. Finally, executable size was not retained as a primary comparison metric because the size of the complete linked executable includes fixed overhead that may obscure differences attributable to the optimized code.

Future work should evaluate more benchmarks and larger search budgets, investigate sequences applied after \texttt{-O3}, expand the available pass set, and provide equivalent initialization information to competing methods. Further experiments should control CPU frequency and affinity, measure code size using the relevant executable sections, and investigate the relationship between surrogate predictions and verified performance.

\section{Conclusion}
This work investigated sequence-aware Bayesian optimization for LLVM compiler pass ordering and compared it with Random Search, a Genetic Algorithm, Tree-structured Parzen Estimator optimization, and a variant without warm start. The evaluation covered five PolyBench/C benchmarks with a budget of 20 evaluations per search and five random seeds. After correctness checking and re-timing, the proposed method achieved a geometric-mean speedup of 0.942 relative to LLVM \texttt{-O3}, compared with 0.952 for TPE, 0.947 for Random Search, and 0.917 for the Genetic Algorithm. No method achieved a verified geometric-mean speedup above \texttt{-O3}, and paired comparisons did not establish a statistically significant verified-speedup advantage for the proposed method. Removing the warm start did not materially change the conclusion. The proposed method showed lower tuning time and fewer failed evaluations in the reported experiment, but these findings concern search behavior rather than a demonstrated improvement in final program performance. Larger budgets, more benchmarks, and further controlled experiments are required to determine its practical value.

\section*{Code Availability}
The source code and experiment scripts for this study are maintained at \url{https://github.com/shauryaraghav/compiler-pass-ordering/}. The results reported in this manuscript correspond to the experimental protocol described in the paper.
